\section{Dirichlet's Theorem on Arithmetic Progressions}
\subsection{Proof of the theorem}
\begin{theorem}[Dirichlet]\label{thm:theorem_dirichlet}
    Let $a,q \in \mathbb{N}$ with $(a,q)=1$. Then, for $x \ge 2$,
    \begin{equation*}
        \sum_{\substack{p\le x \\ p \equiv a (q)}}\frac{1}{p}=  \frac{\log_2 x}{\phi(q)} + O(1),
    \end{equation*}
    where the constant in $O(1)$ depends on $q$.
\end{theorem}
\begin{lemma}\label{lemma:lemma_dirichlet_theorem}
    Let $x\ge2$, $\sigma = 1 + \frac{1}{\log x}$ and $f:\mathbb{N} \to \mathbb{C}$ such that $\lvert f(n)\rvert\le 1$. Then
    \begin{equation*}
        \sum_{p} \frac{f(p)}{p^\sigma} = \sum_{p \le x} \frac{f(p)}{p} + O(1),
    \end{equation*}
    where the constant in $O(1)$ does not depend on $f$ or $x$.
\end{lemma}
\begin{proof}
    We write
    \begin{equation*}
        \sum_{p}\frac{f(p)}{p^\sigma} - \sum_{p\le x}\frac{f(p)}{p} = \sum_{p> x}\frac{f(p)}{p^\sigma} + \sum_{p\le x}f(p)\left(\frac{1}{p^\sigma} - \frac1p\right).
    \end{equation*}
    \emph{The tail.} Since $\lvert f(p)\rvert\le 1$, it is enough to bound $\sum_{p>x}p^{-\sigma}$. We use \cref{thm:Abel_Summation_formula} with $a_n = \mathds{1}_{n \text{ prime}}$, $A = \pi$ and the function $t\mapsto t^{-\sigma}$ on $(x,N]$:
    \begin{equation*}
        \sum_{x<p\le N}\frac{1}{p^\sigma} = \frac{\pi(N)}{N^\sigma} - \frac{\pi(x)}{x^\sigma} + \sigma\int_x^N\frac{\pi(t)}{t^{\sigma+1}}\,dt \le N^{1-\sigma} + \sigma\int_x^\infty\frac{\pi(t)}{t^{\sigma+1}}\,dt.
    \end{equation*}
    Letting $N\to\infty$ and using the Chebyshëv bound $\pi(t)\le K\frac{t}{\log t}$ (\cref{thm:chebyshev}) and $\log t\ge\log x$ for $t \ge x$, we get
    \begin{equation*}
        \sum_{p>x}\frac{1}{p^\sigma}\le K\sigma\int_x^\infty\frac{dt}{t^\sigma\log t}\le\frac{K\sigma}{\log x}\int_x^\infty\frac{dt}{t^\sigma} = \frac{K\sigma}{\log x}\cdot\frac{x^{1-\sigma}}{\sigma-1} = K\sigma e^{-1}\le 2K,
    \end{equation*}
    because $(\sigma-1)\log x = 1$, $x^{1-\sigma} = e^{-1}$ and $\sigma \le 1 + \frac1{\log 2}<3$.

    \emph{The main part.} Using $\lvert f(p)\rvert\le 1$ and the inequality $1-e^{-u} \le u$ for $u \ge 0$, we have
    \begin{equation*}
        \left\lvert \sum_{p \le x} f(p)\left(\frac{1}{p^\sigma} - \frac1p\right) \right\rvert \le\sum_{p \le x} \frac{1-p^{1-\sigma}}{p} = \sum_{ p \le x}\frac{1 - e^{-\frac{\log p}{\log x}}}{p} \le \frac{1}{\log x}\sum_{ p \le x}\frac{\log p}{p }.
    \end{equation*}
    By the first Mertens' estimate (\cref{thm:Mertens} \ref{thm:Mertens:i}), this is $\frac{1}{\log x }\left(\log x + O(1)\right) = O(1)$.
\end{proof}

\begin{proof}[Proof of \Cref{thm:theorem_dirichlet}]
    \emph{Step 1: non-principal characters.} Let $\psi$ be a non-principal character $\pmod q$. We show that
    \begin{equation}\label{eq:sum_non_principal_bound}
        \sum_{p\le x}\frac{\psi(p)}{p} = O(1).
    \end{equation}
    For $\sigma>1$, let $\ell(\sigma,\psi) = \sum_p\sum_{k\ge 1}\frac{\psi(p)^k}{kp^{k\sigma}}$ be the logarithm of $L(\sigma,\psi)$ introduced in the proof of \cref{thm:L-function_neq_0_non-trivial}. The terms with $k\ge2$ are bounded independently of $\sigma$:
    \begin{equation*}
        \left\lvert\sum_p\sum_{k\ge2}\frac{\psi(p)^k}{kp^{k\sigma}}\right\rvert\le\sum_p\sum_{k\ge2}\frac{1}{p^k} = \sum_p\frac{1}{p(p-1)}\le 1,
    \end{equation*}
    so that
    \begin{equation}\label{eq:sum_psi_sigma}
        \sum_p\frac{\psi(p)}{p^\sigma} = \ell(\sigma,\psi) + O(1), \qquad \sigma>1.
    \end{equation}
    We now show that $\ell(\sigma,\psi)$ is bounded for $1<\sigma\le2$. By \cref{thm:analytic_continuation_L-function}, $\sigma\mapsto L(\sigma,\psi)$ is continuous on $[1,2]$. It does not vanish on $(1,2]$ since $L(\sigma,\psi) = e^{\ell(\sigma,\psi)}$, and it does not vanish at $1$ by \cref{thm:L-function_neq_0_non-trivial}. A continuous function $[1,2]\to\mathbb{C}\setminus\{0\}$ has a continuous logarithm $\lambda$ on $[1,2]$, which is bounded since $[1,2]$ is compact. On $(1,2]$, $\ell(\cdot,\psi)$ is continuous (the series converges uniformly on $[1+\delta,2]$ for every $\delta>0$) and $e^{\ell} = e^{\lambda}$, so $\ell-\lambda$ is a continuous function with values in $2\pi i\mathbb{Z}$ on the interval $(1,2]$: it is constant. Hence $\ell(\sigma,\psi)$ is bounded on $(1,2]$.

    Now let $x\ge e$ and $\sigma = 1 + \frac{1}{\log x}\in(1,2]$. By \cref{lemma:lemma_dirichlet_theorem} (with $f = \psi$, since $\lvert\psi(p)\rvert \le 1$) and \cref{eq:sum_psi_sigma},
    \begin{equation*}
        \sum_{p\le x}\frac{\psi(p)}{p} = \sum_p\frac{\psi(p)}{p^\sigma} + O(1) = \ell(\sigma,\psi) + O(1) = O(1).
    \end{equation*}
    For $2\le x<e$, the sum is trivially bounded. This proves \cref{eq:sum_non_principal_bound}.

    \emph{Step 2: detecting the progression.} Since $(a,q)=  1$, the orthogonality relation (\cref{eq:orthogonality_rel_special_case} with $\alpha = p$ and $\beta = a$) gives
    \begin{equation*}
        \sum_{\substack{p \le x \\ p \equiv a (q)}} \frac{1}{p} = \sum_{p\le x}\frac1p\cdot\frac{1}{\phi(q)} \sum_{\chi \bmod q} \overline\chi(a)\chi(p) = \frac{1}{\phi(q)} \sum_{\chi \bmod q}\overline\chi(a) \sum_{p\le x}\frac{\chi(p)}{p}.
    \end{equation*}
    The principal character is like a thorn in our foot, since most of our results don't apply when $\chi = \chi_0$. To manage this situation, we isolate the principal character in our sum:
    \begin{equation*}
        \sum_{\substack{p \le x \\ p \equiv a (q)}} \frac{1}{p} = \frac{1}{\phi(q)} \overline\chi_0(a) \sum_{p\le x}\frac{\chi_0(p)}{p} + \frac{1}{\phi(q)} \sum_{\chi \neq \chi_0}\overline\chi(a) \sum_{p\le x}\frac{\chi(p)}{p}.
    \end{equation*}
    By \cref{eq:sum_non_principal_bound} and $\lvert\overline\chi(a)\rvert = 1$, the second term is a sum of $\phi(q)-1$ bounded terms, so it is $O(1)$. For the first term, $\overline\chi_0(a) = 1$ since $(a,q) = 1$, and $\chi_0(p) = 1$ except for the finitely many primes $p \mid q$. By \cref{thm:Mertens} \ref{thm:Mertens:ii},
    \begin{equation*}
        \sum_{p\le x}\frac{\chi_0(p)}{p} = \sum_{p\le x}\frac1p - \sum_{\substack{p\mid q\\ p \le x}}\frac1p = \log_2 x + O(1).
    \end{equation*}
    We conclude that
    \begin{equation*}
        \sum_{\substack{p \le x \\ p \equiv a (q)}} \frac{1}{p} =\frac{1}{\phi(q)}\log_2 x + O(1).
    \end{equation*}
\end{proof}

\begin{corollary}
    If $(a,d)=1$, then the arithmetic progression $a+nd$ contains infinitely many primes.
\end{corollary}
\begin{proof}
    By \cref{thm:theorem_dirichlet} with $q = d$, $\sum_{p\le x,\ p\equiv a (d)}\frac1p\to\infty$ as $x\to\infty$. If there were only finitely many primes $p \equiv a \pmod d$, this sum would be bounded.
\end{proof}
